
\documentclass{article}
\usepackage{colortbl}
\usepackage{makecell}
\usepackage{multirow}
\usepackage{supertabular}

\begin{document}

\newcounter{utterance}

\centering \large Interaction Transcript for game `cladder', experiment `full\_v1.5\_default', episode 6304 with selfrefine{-}llama{-}3.1{-}8b{-}instruct.
\vspace{24pt}

{ \footnotesize  \setcounter{utterance}{1}
\setlength{\tabcolsep}{0pt}
\begin{supertabular}{c@{$\;$}|p{.15\linewidth}@{}p{.15\linewidth}p{.15\linewidth}p{.15\linewidth}p{.15\linewidth}p{.15\linewidth}}
    \# & \multicolumn{2}{c}{Player} && \multicolumn{2}{c}{Game Master} \\
    \hline

    \theutterance \stepcounter{utterance}  
    & & & \multicolumn{4}{p{0.6\linewidth}}{
        \cellcolor[rgb]{0.9,0.9,0.9}{
            \makecell[{{p{\linewidth}}}]{
                \texttt{\tiny{[P1$\langle$GM]}}
                \texttt{You are an expert in causal inference. The following question is not a typical commonsense query, but rather a meticulously designed question created by a professor specializing in causal inference, intended to assess the students' mastery of the course content.} \\
\texttt{Imagine a self{-}contained, hypothetical world with only the following conditions, and without any unmentioned factors or causal relationships: Zuph has a direct effect on rixq and xevu. Rixq has a direct effect on xevu. The overall probability of rixq is 32\%. The probability of not rixq and xevu is 35\%. The probability of rixq and xevu is 21\%. Is the chance of xevu smaller when observing rixq?} \\
\texttt{Start your answer with 'Yes' or 'No', followed by additional reasoning or evidence to support your explanation.} \\
            }
        }
    }
     \\ \\

    \theutterance \stepcounter{utterance}  
    & \multicolumn{4}{p{0.6\linewidth}}{
        \cellcolor[rgb]{0.561,0.737,0.561}{
            \makecell[{{p{\linewidth}}}]{
                \texttt{\tiny{[P1$\rangle$GM]}}
                \texttt{Yes} \\
\\ 
\texttt{The contradiction in the calculation indicates that our initial assumption that P(xevu and not rixq) + P(xevu and rixq) = P(xevu) is incorrect. However, we can use the fact that P(xevu) = P(xevu and rixq) + P(xevu and not rixq) and P(xevu and not rixq) + P(not rixq and xevu) = P(xevu) to find P(xevu and not rixq).} \\
\\ 
\texttt{Since P(xevu and not rixq) + P(not rixq and xevu) = P(xevu), we can use the given values to find P(xevu and not rixq):} \\
\\ 
\texttt{P(xevu and not rixq) + P(not rixq and xevu) = P(xevu)} \\
\texttt{P(xevu and not rixq) + 35\% = P(xevu)} \\
\\ 
\texttt{However, we also know that P(xevu) = P(xevu and rixq) + P(xevu and not rixq), and P(xevu and rixq) = 21\%. Therefore, we can substitute P(xevu) in the previous equation:} \\
\\ 
\texttt{P(xevu and not rixq) + 35\% = P(xevu and rixq) + P(xevu and not rixq)} \\
\texttt{35\% = 21\%} \\
\\ 
\texttt{This is another contradiction, which means our initial assumption that P(xevu and not rixq) + P(not rixq and xevu) = P(xevu) is incorrect. However, we can use the fact that P(xevu) = P(xevu and rixq) + P(xevu and not rixq) and P(xevu and not rixq) + P(not rixq and xevu) = P(xevu) to find P(xevu and not rixq) by using the given values:} \\
\\ 
\texttt{P(xevu and not rixq) + P(not rixq and xevu) = P(xevu)} \\
\texttt{P(xevu and not rixq) + 35\% = P(xevu and rixq) + P(xevu and not rixq)} \\
\texttt{P(xevu and not rixq) + 35\% = P(xevu and rixq) + P(xevu and not rixq)} \\
\texttt{P(xevu and not rixq) + 35\% = 21\% + P(xevu and not rixq)} \\
\texttt{P(xevu and not rixq) + 35\% = P(xevu and not rixq) + 21\%} \\
\texttt{35\% = 21\%} \\
\\ 
\texttt{However, we can use the fact} \\
            }
        }
    }
    & & \\ \\

    \theutterance \stepcounter{utterance}  
    & & & \multicolumn{2}{p{0.3\linewidth}}{
        \cellcolor[rgb]{0.9,0.9,0.9}{
            \makecell[{{p{\linewidth}}}]{
                \texttt{\tiny{[GM$|$GM]}}
                \texttt{Yes} \\
            }
        }
    }
    & & \\ \\

    \theutterance \stepcounter{utterance}  
    & & & \multicolumn{2}{p{0.3\linewidth}}{
        \cellcolor[rgb]{0.9,0.9,0.9}{
            \makecell[{{p{\linewidth}}}]{
                \texttt{\tiny{[GM$|$GM]}}
                \texttt{no} \\
            }
        }
    }
    & & \\ \\

    \theutterance \stepcounter{utterance}  
    & & & \multicolumn{2}{p{0.3\linewidth}}{
        \cellcolor[rgb]{0.9,0.9,0.9}{
            \makecell[{{p{\linewidth}}}]{
                \texttt{\tiny{[GM$|$GM]}}
                \texttt{game\_result = LOSE} \\
            }
        }
    }
    & & \\ \\

\end{supertabular}
}

\end{document}
